% ACM Multimedia 2026 Submission
% Cultural-Aware Evaluation Routing for AI Art Generation (v2 rewrite)
\documentclass[sigconf,authordraft]{acmart}

% Additional packages (acmart includes booktabs, xcolor, hyperref, etc.)
\usepackage{multirow}
\usepackage{subcaption}
\usepackage{graphicx}
\usepackage{algorithm}
\usepackage{algpseudocode}
\usepackage{tikz}
\usetikzlibrary{shapes.geometric, arrows.meta, positioning, calc, fit, backgrounds}

% Custom commands
\newcommand{\vulca}{\textsc{Vulca}}
\newcommand{\system}{\textsc{Vulca-Agent}}
\newcommand{\Lx}[1]{L#1}
\newcommand{\todo}[1]{\textcolor{red}{[TODO: #1]}}

% ---- ACM metadata ----
\copyrightyear{2026}
\acmYear{2026}
\setcopyright{acmlicensed}
\acmConference[MM '26]{Proceedings of the 34th ACM International Conference on Multimedia}{November 10--14, 2026}{Rio de Janeiro, Brazil}
\acmBooktitle{Proceedings of the 34th ACM International Conference on Multimedia (MM '26), November 10--14, 2026, Rio de Janeiro, Brazil}
\acmDOI{10.1145/nnnnnnn.nnnnnnn}
\acmISBN{978-x-xxxx-xxxx-x/26/11}

\setlength{\emergencystretch}{3em}

\begin{document}

\title{Cultural-Aware Art Generation via Tradition-Specific Evaluation Routing}

\author{Haorui Yu}
\email{yuhaorui48@gmail.com}
\affiliation{%
  \institution{Independent Researcher}
  \country{China}
}

% ---- CCS Concepts ----
\begin{CCSXML}
<ccs2012>
  <concept>
    <concept_id>10010147.10010178.10010224</concept_id>
    <concept_desc>Computing methodologies~Multi-agent systems</concept_desc>
    <concept_significance>500</concept_significance>
  </concept>
  <concept>
    <concept_id>10010147.10010178.10010179</concept_id>
    <concept_desc>Computing methodologies~Natural language processing</concept_desc>
    <concept_significance>300</concept_significance>
  </concept>
  <concept>
    <concept_id>10010405.10010469</concept_id>
    <concept_desc>Applied computing~Arts and humanities</concept_desc>
    <concept_significance>500</concept_significance>
  </concept>
</ccs2012>
\end{CCSXML}

\ccsdesc[500]{Computing methodologies~Multi-agent systems}
\ccsdesc[300]{Computing methodologies~Natural language processing}
\ccsdesc[500]{Applied computing~Arts and humanities}

\keywords{cultural AI, multi-agent systems, evaluation routing, art generation, dimension weighting}

% ============================================================
\begin{abstract}
% ============================================================

Current text-to-image models generate culturally decontextualized art: they can produce visually appealing images but routinely violate tradition-specific conventions---rendering \emph{xieyi} ink paintings with photorealistic textures, or Islamic geometric patterns with figurative elements.
We present \system{}, a multi-agent evaluation pipeline that grounds art generation in cultural knowledge through five cognitive layers (\Lx{1}--\Lx{5}): Visual Perception, Technical Analysis, Cultural Context, Critical Interpretation, and Philosophical Aesthetics.
\system{} implements a \emph{Scout\,$\to$\,Draft\,$\to$\,Critic\,$\to$\,Queen} pipeline where:
(1)~a \textbf{Scout} agent retrieves cultural evidence via semantic search over a curated knowledge base;
(2)~a \textbf{Draft} agent generates candidate images grounded in retrieved evidence;
(3)~a \textbf{Critic} agent evaluates candidates across \Lx{1}--\Lx{5} using a hybrid rule+LLM approach with \emph{tradition-specific dimension weighting}; and
(4)~a \textbf{Queen} agent decides whether to accept, rerun, or request additional evidence.
The key innovation is a \emph{cultural pipeline router} that selects tradition-specific dimension weights across 9 art traditions (e.g., amplifying \Lx{5} Philosophical Aesthetics for Chinese xieyi, \Lx{2} Technical Analysis for Islamic geometric art), enabling the evaluation to reflect each tradition's priorities rather than applying uniform criteria.
Across 480 real-mode runs (16 conditions $\times$ 30 culturally diverse tasks spanning four experimental phases), we show that cultural routing alone improves mean weighted score from $0.810$ to $0.846$ ($\Delta = +0.036$, $p < 0.001$, Cohen's $d = 0.95$), accounting for 95\% of the full pipeline's effect ($+0.038$, $d = 0.94$).
Gains concentrate on \Lx{5} Philosophical Aesthetics ($+0.134$, $p < 0.001$) and \Lx{3} Cultural Context ($+0.056$, $p < 0.001$).
Supplementary experiments reveal three complementary findings: (1)~multi-round iteration adds only $+0.002$ ($p = 0.35$) under CLIP-based evaluation, identifying scoring granularity as the bottleneck; (2)~replacing CLIP with VLM-based per-layer scoring resolves this bottleneck ($\Delta = +0.077$, $d = 2.10$, $p < 0.001$), while revealing that LLM and VLM evaluation functions overlap---making FLUX\,+\,VLM routing without LLM critic the strongest single configuration ($0.887$); and (3)~replacing the FLUX diffusion backend with a hybrid AR+Diffusion generator (NB2) further raises the peak to $\mathbf{0.917}$ ($+0.107$ vs.\ baseline, $d = 2.97$), but the multi-round agent loop yields only $+0.002$ ($p = 0.70$)---confirming that self-correction has \emph{diminishing returns} as generator quality improves.

\end{abstract}

\maketitle

% ============================================================
\section{Introduction}
\label{sec:intro}
% ============================================================

The rapid advancement of text-to-image generation models---from Stable Diffusion~\cite{rombach2022high} to FLUX~\cite{flux2024} and DALL-E~3~\cite{dalle3}---has democratized visual content creation.
However, these models treat cultural art traditions as mere stylistic filters rather than rich semiotic systems with specific conventions, taboos, and philosophical underpinnings.
When prompted to generate ``Chinese xieyi ink painting of bamboo,'' models like Stable Diffusion frequently produce photorealistic bamboo with ink wash effects---violating the fundamental xieyi principle of \emph{yibi} (意筆, meaning-laden brushwork) where each stroke carries philosophical weight beyond visual representation.

This cultural decontextualization is not merely an aesthetic shortcoming but a \emph{categorically different} problem from generic image quality.
A tradition like Islamic geometric art encodes theological principles in its tessellation patterns: the infinite extension of arabesques symbolizes the boundlessness of creation, while the prohibition of figurative representation reflects aniconism rooted in centuries of scholarly discourse.
Generating ``an Islamic geometric pattern'' with a human face in the corner is not a quality failure---it is a \emph{semantic} failure that reveals the absence of cultural reasoning entirely.
Similarly, Chinese xieyi painting privileges expressive emptiness (\emph{liu bai}/留白) as a philosophical statement, not merely a composition choice; photorealistic gap-filling violates the art form's core epistemology.

Current models cannot make these distinctions because they optimize for visual fidelity, not cultural coherence.
Recent work confirms this is not a gap that can be closed by scaling models or refining prompts alone.
Culture-TRIP~\cite{jeong2025culturetrip} demonstrates that standard text-to-image models fail to capture culturally specific concepts even with carefully engineered prompts, requiring iterative refinement grounded in external cultural knowledge.
Broader surveys document persistent geo-cultural biases that resist dataset-level mitigation~\cite{wan2024survey}, while multilingual evaluation reveals that models collapse diverse cultural prompts into culturally neutral outputs~\cite{shi2025culture}.
These findings point to a structural limitation: generation models lack the representational capacity to simultaneously optimize for visual fidelity and cultural coherence.
The remaining alternative is an \emph{external} evaluation pipeline that assesses generated output against tradition-specific cultural criteria---an architectural intervention rather than a parametric one.

In our 30-task evaluation, the single-pass FLUX baseline achieves \Lx{5} Philosophical Aesthetics scores averaging only 0.705 and \Lx{3} Cultural Context scores of 0.816---visually acceptable but culturally shallow.
The gap lies in \emph{cognitive depth}: a human art critic evaluates paintings through layered processes---from basic visual perception through technical analysis, cultural contextualization, critical interpretation, to philosophical appreciation~\cite{panofsky1939studies, barrett2011criticizing}.
These are not independent checks but a \emph{dependency chain}: cultural judgment (\Lx{3}) presupposes correct visual and technical perception (\Lx{1}--\Lx{2}), and philosophical evaluation (\Lx{5}) requires grounding in cultural context (\Lx{3}--\Lx{4}).
No existing generation system replicates this layered cognitive process, let alone uses it to \emph{drive} iterative improvement of generated outputs.

We address this gap with \system{}, a multi-agent framework built on the \vulca{} five-layer evaluation hierarchy (\Lx{1}--\Lx{5})~\cite{yu2025vulca, yu2025vulcabench}.
Our system makes three contributions:

\begin{enumerate}
    \item \textbf{Tradition-Specific Evaluation Routing.} A cultural pipeline router selects tradition-appropriate dimension weights across 9 art traditions (e.g., amplifying \Lx{5} for Chinese xieyi, \Lx{2} for Islamic geometric art), enabling culturally differentiated evaluation rather than uniform criteria. This single mechanism accounts for 95\% of observed quality gains.

    \item \textbf{Structured Cultural Grounding.} A Scout agent retrieves tradition-specific evidence (terminology, exemplar artworks, taboos) via semantic search and packages it as an \emph{EvidencePack} that constrains downstream generation and evaluation.

    \item \textbf{Empirical Decomposition of Cultural Improvement.} Controlled ablation across 480 real-mode runs (16 conditions $\times$ 30 tasks) spanning four experimental phases. We demonstrate that: cultural routing alone yields $+0.036$ ($p < 0.001$, $d = 0.95$); multi-round iteration under CLIP-based scoring adds only $+0.002$ ($p = 0.35$), identifying scoring granularity as the bottleneck under CLIP evaluation; VLM-based per-layer scoring resolves this bottleneck ($+0.077$, $d = 2.10$), while showing that LLM critics become redundant under high-fidelity VLM evaluation; and generator upgrade (FLUX$\to$NB2) raises the peak to $0.917$ while confirming that self-correction loops have diminishing returns as generator quality improves.
\end{enumerate}

% ============================================================
\section{Related Work}
\label{sec:related}
% ============================================================

\paragraph{Cultural AI Art Evaluation.}
The \vulca{} framework~\cite{yu2025vulca} introduced a five-dimensional evaluation protocol for AI-generated cultural art, establishing the \Lx{1}--\Lx{5} hierarchy.
\vulca{}-Bench~\cite{yu2025vulcabench} operationalized this into a benchmark of 7,410 samples across 8 cultural traditions.
Other evaluation frameworks include ArtCritic~\cite{yu2025artcritique} for cross-cultural critique assessment and GalleryGPT~\cite{gallerygpt2024} for visual art understanding.
Our work extends these evaluation frameworks into an \emph{active evaluation pipeline} where tradition-specific routing drives cultural quality improvement.

\paragraph{Multi-Agent Systems for Generation.}
Multi-agent architectures have shown promise in code generation~\cite{hong2024metagpt}, scientific discovery~\cite{boiko2023autonomous}, and creative writing~\cite{wang2024chatdev}.
In the visual domain, RPG-Diffusion\-Master~\cite{yang2024mastering} uses LLMs to decompose complex prompts for diffusion models.
However, none of these systems incorporate cultural knowledge bases or tradition-specific evaluation routing grounded in domain-specific criteria.

\paragraph{Self-Correction in LLMs and Diffusion Models.}
Self-refinement~\cite{madaan2024self} showed that LLMs iteratively improve outputs via self-feedback.
In image generation, ControlNet~\cite{zhang2023adding} enables spatial control and IP-Adapter~\cite{ye2023ip} provides style transfer.
Our FixItPlan protocol combines these technical tools with cultural diagnostic information to enable \emph{targeted} self-correction at specific cognitive layers.

\paragraph{VLM-Based Evaluation Metrics.}
TIFA~\cite{hu2023tifa} measures text-image faithfulness via VQA; VIEScore~\cite{ku2024viescore} uses MLLMs for explainable image scoring; MLLM-as-a-Judge~\cite{chen2024mllmjudge} benchmarks multimodal judgment, confirming pairwise comparison reliable but scoring granularity challenging.
Our VLMCritic applies per-layer VLM evaluation across \Lx{1}--\Lx{5} to resolve the granularity bottleneck that CLIP-based scoring cannot overcome (\S\ref{sec:analysis}).

% ============================================================
\section{The \Lx{1}--\Lx{5} Framework}
\label{sec:framework}
% ============================================================

Our evaluation hierarchy synthesizes two foundational traditions in art criticism.
Panofsky's iconographic method~\cite{panofsky1939studies} distinguishes three strata of meaning: \emph{pre-iconographic} description (what is depicted), \emph{iconographic} analysis (what it symbolizes within a tradition), and \emph{iconological} interpretation (what it reveals about a culture's worldview).
Barrett's criticism pedagogy~\cite{barrett2011criticizing} operationalizes this into four stages: describing, analyzing, interpreting, and evaluating.
The \vulca{} framework~\cite{yu2025vulcabench} extends these into a five-layer hierarchy designed to be \emph{computationally actionable}---each layer maps to specific measurable criteria rather than open-ended judgment:

\noindent
\textbf{\Lx{1} Visual Perception}: what is seen---composition, color palette, spatial arrangement. Corresponds to Panofsky's pre-iconographic level and is evaluable by deterministic rules.
\textbf{\Lx{2} Technical Analysis}: how it is made---medium-specific techniques (ink saturation in sumi-e, tessellation precision in Islamic geometry). Still largely rule-assessable.
\textbf{\Lx{3} Cultural Context}: what it means within a tradition---adherence to conventions, symbolic systems, and taboo avoidance. This is the first layer requiring external cultural knowledge, marking the boundary between ``visual quality'' and ``cultural quality.''
\textbf{\Lx{4} Critical Interpretation}: deeper engagement with art historical discourse---understanding an artwork's position within a tradition's evolution.
\textbf{\Lx{5} Philosophical Aesthetics}: transcendent value---whether the work achieves qualities like wabi-sabi impermanence, sublime terror, or xieyi's ``meaning beyond form'' (意在筆先). The most difficult layer to evaluate and the one where current AI systems fail most dramatically.

\noindent
Critically, these layers form a \textbf{dependency chain} rather than independent dimensions: accurate \Lx{3} cultural judgment requires correct \Lx{1}--\Lx{2} perception (one cannot identify a taboo violation without first recognizing what is depicted), and meaningful \Lx{5} evaluation presupposes \Lx{3}--\Lx{4} understanding (philosophical depth without cultural grounding is vacuous).
This dependency has two direct engineering consequences: (1)~our Critic evaluates bottom-up, so that lower-layer failures are caught before higher-layer evaluation is attempted; and (2)~our repair strategy targets the \emph{lowest failing layer}, since fixing a \Lx{2} technical deficiency may cascade upward to improve \Lx{3}--\Lx{5} scores without additional intervention.

% ============================================================
\section{System Architecture}
\label{sec:system}
% ============================================================

Given that prompt engineering alone cannot enforce structural cultural constraints~\cite{jeong2025culturetrip, shi2025culture} (\S\ref{sec:intro}), the remaining path is an external evaluation pipeline.
But why four agents rather than, say, a single LLM that retrieves, generates, evaluates, and decides?
The answer derives from the failure mode that aggregate evaluation enables: when a single system evaluates its own output, it can trade off between dimensions---sacrificing cultural accuracy to maintain visual quality, or suppressing lower-level detail to inflate higher-level scores.
Our four-agent design prevents this by separating concerns so that each agent is accountable for a distinct cognitive function:

\begin{itemize}
    \item \textbf{Scout} $\to$ \emph{cultural knowledge that generators lack}. Diffusion models encode statistical co-occurrences between tokens and pixels, not the recursive symbolic references that cultural traditions carry. The Scout makes this external knowledge explicit via the EvidencePack.
    \item \textbf{Draft} $\to$ \emph{generation as a black box}. We accept the generator's opacity and focus on constraining its input (via evidence-derived prompts) and filtering its output (via the Critic).
    \item \textbf{Critic} $\to$ \emph{independent, per-layer evaluation}. Evaluating each \Lx{1}--\Lx{5} layer separately prevents aggregate metrics from masking dimension-level failures---a pattern analogous to reward hacking in RL, where Gap reduction is achieved via lower-level suppression rather than genuine improvement.
    \item \textbf{Queen} $\to$ \emph{cross-dimension protection}. The Queen tracks confirmed vs.\ pending dimensions, ensuring that improvements in one layer do not degrade another---directly addressing the observation that unconstrained optimization sacrifices the most culturally vulnerable dimensions first.
\end{itemize}

\noindent
The separation is not merely architectural convenience---it reflects the observation that \emph{retrieval}, \emph{generation}, \emph{evaluation}, and \emph{decision-making} are fundamentally different cognitive tasks that benefit from specialized processing (Figure~\ref{fig:architecture}).

\begin{figure*}[t]
    \centering
    % VULCA Agent Loop Architecture Diagram (inline version for main.tex)
% Requires in preamble: \usepackage{tikz}
% \usetikzlibrary{shapes.geometric, arrows.meta, positioning, calc, fit, backgrounds}

\definecolor{scoutC}{HTML}{3B82F6}
\definecolor{draftC}{HTML}{10B981}
\definecolor{criticC}{HTML}{F59E0B}
\definecolor{queenC}{HTML}{8B5CF6}
\definecolor{routerC}{HTML}{EC4899}
\definecolor{loopC}{HTML}{EF4444}
\definecolor{vlmC}{HTML}{06B6D4}
\definecolor{bdr}{HTML}{94A3B8}

\resizebox{\textwidth}{!}{%
\begin{tikzpicture}[
    >=Stealth,
    agent/.style={
        rectangle, rounded corners=5pt, draw=#1, fill=#1!8,
        line width=1pt, minimum width=2.2cm, minimum height=0.9cm,
        font=\sffamily\footnotesize\bfseries, align=center, text=#1!85!black
    },
    sub/.style={font=\sffamily\tiny, text=#1!60!black},
    data/.style={
        rectangle, rounded corners=2pt, draw=bdr, fill=white,
        line width=0.5pt, font=\sffamily\tiny, align=center,
        minimum height=0.4cm, inner sep=2pt
    },
    lbox/.style={
        rectangle, rounded corners=1.5pt, draw=#1!50, fill=#1!10,
        line width=0.4pt, font=\sffamily\tiny\bfseries, text=#1!70!black,
        minimum width=0.48cm, minimum height=0.33cm, inner sep=1pt
    },
    arr/.style={->, line width=0.7pt, color=#1!55!black},
    darr/.style={->, line width=0.7pt, dashed, color=#1!55!black},
    note/.style={font=\sffamily\tiny\itshape, text=gray!50!black},
]

% ==================== ROW 1: Input → Scout → Draft → Cands ====================
\node[data, minimum width=1.4cm] (input) at (0,0) {Task $+$ Tradition};

\node[agent=scoutC] (scout) at (2.7, 0) {Scout};
\node[sub=scoutC, below=0.01cm of scout] {Evidence Retrieval};

\node[data, below=0.55cm of scout, xshift=-0.7cm] (faiss) {FAISS};
\node[data, below=0.55cm of scout] (terms) {Terms};
\node[data, below=0.55cm of scout, xshift=0.7cm] (taboo) {Taboo};

\node[data, fill=scoutC!5, minimum width=1.4cm] (epack) at (5.1, 0) {EvidencePack};

\node[agent=draftC] (draft) at (7.3, 0) {Draft};
\node[sub=draftC, below=0.01cm of draft] {Image Generation};

\node[data, below=0.55cm of draft, xshift=-0.9cm] (sd) {SD\,1.5};
\node[data, below=0.55cm of draft, xshift=-0.3cm] (flx) {FLUX};
\node[data, below=0.55cm of draft, xshift=0.3cm] (nb2) {\textbf{NB2}};
\node[data, below=0.55cm of draft, xshift=0.9cm] (cn) {CtrlNet};

\node[data, fill=draftC!5, minimum width=1.2cm] (cands) at (9.5, 0) {$n$ Cands};

% ==================== ROW 2: Queen ← Scores ← Critic ====================
\node[agent=criticC] (critic) at (7.8, -2.8) {Critic};
\node[sub=criticC, below=0.01cm of critic] {L1--L5 Hybrid};

% L1-L5 boxes — Rule (L1-L2), LLM (L3-L5), VLM (L1-L5)
\node[lbox=scoutC] (l1) at (6.55, -3.7) {L1};
\node[lbox=scoutC] (l2) at (7.0, -3.7) {L2};
\node[lbox=criticC] (l3) at (7.5, -3.7) {L3};
\node[lbox=criticC] (l4) at (8.0, -3.7) {L4};
\node[lbox=queenC] (l5) at (8.5, -3.7) {L5};
\draw[draw=bdr!50, line width=0.3pt] ([yshift=-0.15cm]l1.south west) -- ([yshift=-0.15cm]l2.south east);
\node[note] at (6.78, -4.08) {Rule};
\draw[draw=bdr!50, line width=0.3pt] ([yshift=-0.15cm]l3.south west) -- ([yshift=-0.15cm]l5.south east);
\node[note] at (8.0, -4.08) {LLM};

% VLM evaluation path (spans L1-L5)
\node[lbox=vlmC] (vlm1) at (6.55, -4.5) {L1};
\node[lbox=vlmC] (vlm2) at (7.0, -4.5) {L2};
\node[lbox=vlmC] (vlm3) at (7.5, -4.5) {L3};
\node[lbox=vlmC] (vlm4) at (8.0, -4.5) {L4};
\node[lbox=vlmC] (vlm5) at (8.5, -4.5) {L5};
\draw[draw=bdr!50, line width=0.3pt] ([yshift=-0.15cm]vlm1.south west) -- ([yshift=-0.15cm]vlm5.south east);
\node[note, text=vlmC!70!black] at (7.53, -4.88) {\textbf{VLM} (Gemini)};

\node[data, fill=criticC!5, minimum width=1.2cm, minimum height=0.5cm] (scores) at (5.5, -2.8) {Scores $+$\\[-1pt]Signals};

\node[agent=queenC] (queen) at (3.2, -2.8) {Queen};
\node[sub=queenC, below=0.01cm of queen] {Decision Maker};

% Outputs
\node[data, fill=draftC!10, draw=draftC!60, minimum width=1.1cm, font=\sffamily\tiny\bfseries] (accept) at (2.3, -3.9) {Accept};
\node[data, minimum width=1.1cm] (stop) at (4.1, -3.9) {Stop};

% Cultural Router
\node[rectangle, rounded corners=3pt, draw=routerC, fill=routerC!6,
      line width=0.7pt, minimum width=2.0cm, minimum height=0.55cm,
      font=\sffamily\footnotesize\bfseries, text=routerC!80!black] (router) at (10.5, -1.7) {Cultural Router};
\node[note, below=0.01cm of router] {9 traditions $\to$ L-weights};

% ==================== MAIN FLOW ====================
\draw[arr=gray] (input) -- (scout);
\draw[arr=scoutC] (scout) -- (epack);
\draw[arr=scoutC] (epack) -- (draft);
\draw[arr=draftC] (draft) -- (cands);

% U-turn: Candidates → Critic
\draw[arr=draftC] (cands.south) |- (critic.east);

% Critic → Scores → Queen
\draw[arr=criticC] (critic) -- (scores);
\draw[arr=criticC] (scores) -- (queen);

% Queen → Accept / Stop
\draw[arr=draftC] (queen.south) -- ++(0, -0.2) -| (accept.north);
\draw[arr=gray] (queen.south) -- ++(0, -0.2) -| (stop.north);

% Router → Critic (route below to avoid Cands→Critic vertical overlap at x=9.5)
\draw[darr=routerC] (router.south) -- ++(0, -0.8) -| (critic.north east);

% Tradition → Router (route along bottom to avoid crossing central flow)
\draw[darr=routerC] (input.south) -- ++(0, -5.3) -| (router.south);

% ==================== SELF-CORRECTION LOOPS ====================

% Loop 1: Rerun + FixItPlan (Queen → Draft, going above)
\coordinate (loop1top) at (5.2, 1.2);
\draw[->, line width=0.8pt, dashed, color=loopC!65]
    (queen.north) -- (queen.north |- loop1top) -- (draft.north |- loop1top)
    node[note, midway, above, text=loopC!70, font=\sffamily\tiny\bfseries\itshape] {Rerun $+$ FixItPlan}
    -- (draft.north);

% Loop 2: NeedMoreEvidence (Queen → Scout, going further above)
\coordinate (loop2top) at (3.0, 1.7);
\draw[->, line width=0.65pt, dashed, color=scoutC!50]
    ([xshift=-0.2cm]queen.north) -- ([xshift=-0.2cm]queen.north |- loop2top) -- (scout.north |- loop2top)
    node[note, midway, above, text=scoutC!60, font=\sffamily\tiny\bfseries\itshape] {NeedMoreEvidence}
    -- (scout.north);

% ==================== BACKGROUND ====================
\begin{scope}[on background layer]
    \node[fill=gray!3, rounded corners=8pt, inner sep=8pt,
          fit=(scout)(draft)(critic)(queen)(cands)(epack)(scores)(accept)(stop)(l1)(l5)(vlm1)(vlm5)(router)] {};
\end{scope}

\end{tikzpicture}%
}

    \caption{The \system{} pipeline. The Scout retrieves cultural evidence, the Draft generates candidates grounded in an EvidencePack, the Critic evaluates across \Lx{1}--\Lx{5} using a hybrid rule+LLM approach, and the Queen decides whether to accept, rerun with a FixItPlan, or request more evidence. Dashed arrows show self-correction loops; the Cultural Router adjusts \Lx{1}--\Lx{5} weights per tradition.}
    \label{fig:architecture}
\end{figure*}

% -----------------------------------------------------------
\subsection{Scout: Cultural Evidence Retrieval}
\label{sec:scout}

The Scout agent grounds generation in cultural knowledge through three channels:
(1)~\textbf{Semantic search} via FAISS~\cite{johnson2019billion} with MiniLM-L6-v2 and CLIP~\cite{radford2021learning} retrieves relevant samples and visual references from \vulca{}-Bench;
(2)~\textbf{Terminology lookup} from a curated dictionary of 200+ cultural terms with L-level associations (e.g., \emph{cun fa}/皴法 tagged as \Lx{2}-relevant);
(3)~\textbf{Taboo detection} via rule-based scanning for cultural violations.
The Scout packages findings as an \textbf{EvidencePack}---a structured protocol containing anchor terms, exemplar references, composition constraints, style directives, and taboo warnings.

% -----------------------------------------------------------
\subsection{Draft: Evidence-Grounded Generation}
\label{sec:draft}

The Draft agent converts the EvidencePack into a generation prompt and produces $n$ candidate images (default $n{=}4$) using a configurable image generation backend:

Supported backends include Stable Diffusion~1.5~\cite{rombach2022high} (local GPU, fp16), FLUX.1-schnell~\cite{flux2024} (API-based), NB2 (Gemini image generation with internal thinking~\cite{team2024gemini}, AR+Diffusion hybrid), ControlNet~\cite{zhang2023adding} (structural inpainting), and FLUX Fill (diffusion-based masked inpainting for Rerun-Local mode via FixItPlan).
Each candidate uses a different random seed while sharing the same evidence-derived prompt, enabling diversity within cultural constraints.

% -----------------------------------------------------------
\subsection{Critic: Hybrid \Lx{1}--\Lx{5} Evaluation}
\label{sec:critic}

The Critic agent evaluates each candidate across all five layers using a hybrid approach:

\paragraph{Rule-Based Scoring (\Lx{1}--\Lx{2}).}
Deterministic rules provide reproducible baseline scores.
\Lx{1} awards bonuses for style keyword matching, terminology coverage, and prompt specificity;
\Lx{2} rewards appropriate sampling parameters, model capability, and evidence utilization.

\paragraph{LLM-Based Scoring (\Lx{3}--\Lx{5}).}
For layers requiring cultural judgment, we employ a two-phase ReAct pattern with a language model.
In the \emph{exploration} phase, the LLM reasons over the \emph{EvidencePack}---terminology, tradition labels, taboo flags, evidence coverage---to score cultural adherence, interpretive depth, and aesthetic achievement; this is a \emph{text-reasoning} process grounded in retrieved evidence, not direct image analysis.
In the \emph{forced submission} phase, a fresh context forces structured scoring output.
Scores are merged: $s_{\text{merged}} = 0.3 \cdot s_{\text{rule}} + 0.7 \cdot s_{\text{llm}}$ for activated layers.
A v3 extension replaces CLIP-based image similarity with a \textbf{VLM Critic} (Gemini) that directly evaluates generated images across \Lx{1}--\Lx{5} from visual input---a distinct evaluation path; ablation results (\S\ref{sec:analysis}) show VLM outperforms CLIP by $+0.041$ and shows substantial functional overlap with the text-reasoning LLM critic.

\paragraph{Cross-Layer Signals and Dynamic Weights.}
Five diagnostic rules detect inter-layer inconsistencies (e.g., $|\Lx{5} - \Lx{1}| > 0.3$ triggers \textsc{Reinterpret}; $\Lx{3} < 0.5$ triggers \textsc{Cultural\_Violation}; taboo detection triggers immediate \textsc{Taboo\_Override}).
These signals propagate to the Queen for decision-making and modulate layer weights dynamically through confidence adjustment, round decay, signal boosting, and deviation capping.

% -----------------------------------------------------------
\subsection{Queen: Decision Orchestration}
\label{sec:queen}

The Queen agent receives the best candidate's weighted score and cross-layer signals, then decides:

\begin{itemize}
    \item \textbf{Accept}: $w_{\text{total}} \geq \theta_{\text{accept}}$ and no critical cross-layer signals.
    \item \textbf{Rerun}: $w_{\text{total}} < \theta_{\text{accept}}$ or cross-layer signals demand correction.
    \item \textbf{Rerun-Local}: Like rerun, but with a \emph{FixItPlan} that specifies spatial masking and prompt modifications for targeted inpainting via ControlNet.
    \item \textbf{Stop}: Maximum rounds reached or diminishing returns detected.
\end{itemize}

An early-stop threshold ($\theta_{\text{early}} = 0.93$) bypasses the full evaluation pipeline for near-perfect candidates.

% -----------------------------------------------------------
\subsection{FixItPlan: Targeted Self-Correction}
\label{sec:fixitplan}

When the Critic identifies layers scoring below threshold, it generates a \textbf{FixItPlan} specifying target layers, prompt modifications (e.g., ``add \emph{cun fa} texture strokes''), and a spatial mask strategy for inpainting.
For $<$3 failing layers the Queen activates targeted inpainting (ControlNet/FLUX Fill) preserving successful regions; for $\geq$3 failures, full regeneration is triggered.

% -----------------------------------------------------------
\subsection{Cultural Pipeline Router}
\label{sec:router}

Different art traditions prioritize different cognitive layers.
We define tradition-specific weight tables for 9 traditions:

\begin{table}[t]
\centering
\small
\begin{tabular}{@{}lccccc@{}}
\toprule
\textbf{Tradition} & \Lx{1} & \Lx{2} & \Lx{3} & \Lx{4} & \Lx{5} \\
\midrule
Chinese Xieyi   & .10 & .15 & .20 & .25 & \textbf{.30} \\
Chinese Gongbi  & .15 & \textbf{.30} & .25 & .15 & .15 \\
Japanese        & .15 & .20 & .20 & .20 & \textbf{.25} \\
Islamic         & .25 & \textbf{.30} & .15 & .15 & .15 \\
African         & .15 & .15 & \textbf{.30} & .20 & .20 \\
Western Acad.   & .15 & .20 & .15 & \textbf{.25} & .25 \\
S.\ Asian       & .15 & .15 & .20 & .25 & \textbf{.25} \\
Watercolor      & .20 & .25 & .15 & .20 & .20 \\
\bottomrule
\end{tabular}
\caption{\Lx{1}--\Lx{5} weight allocation by tradition (8 of 9 shown; ``default'' omitted). Bold indicates the dominant layer. All rows sum to 1.0.}
\label{tab:cultural_weights}
\end{table}

The router also selects pipeline variants: \emph{chinese\_xieyi} uses atomic (single-pass) generation honoring the yiqi (一氣) principle; \emph{western\_academic} uses three-step progressive refinement.

% ============================================================
\section{Experimental Setup}
\label{sec:experiments}
% ============================================================

\subsection{Task Suite}

We use 30 culturally diverse generation tasks drawn from \vulca{}-Bench~\cite{yu2025vulcabench}:
\begin{itemize}
    \item \textbf{Poetic} (10): Tasks requiring interpretation of literary imagery (e.g., ``Bada Shanren style of two fish expressing cosmic solitude'').
    \item \textbf{Cultural} (10): Tasks testing tradition-specific conventions (e.g., ``Persian miniature depicting a royal court scene with gold leaf'').
    \item \textbf{Taboo} (10): Tasks involving culturally sensitive content requiring careful handling (e.g., ``Chinese dragon in a politically sensitive context'').
\end{itemize}

\subsection{Ablation Conditions}

We design 16 conditions across four experimental phases to isolate the contribution of each system component (Table~\ref{tab:conditions}):

\begin{table}[t]
\centering
\small
\begin{tabular}{@{}cllccc@{}}
\toprule
 & \textbf{Model} & \textbf{Critic} & \textbf{Loop} & \textbf{Route} & $\boldsymbol{\theta}$ \\
\midrule
A & SD1.5 & Rule & $\times$ & $\times$ & .80 \\
B & SD1.5 & Rule+LLM & $\checkmark$ & $\checkmark$ & .80 \\
C & FLUX & Rule & $\times$ & $\times$ & .80 \\
\textbf{D} & \textbf{FLUX} & \textbf{Rule+LLM} & $\checkmark$ & $\checkmark$ & .80 \\
E & FLUX & Rule & $\checkmark$ & $\checkmark$ & .80 \\
F & FLUX & Rule+LLM & $\times$ & $\checkmark$ & .80 \\
\midrule
\multicolumn{6}{@{}l}{\emph{v2 supplementary (tightened thresholds):}} \\
C\textsuperscript{+} & FLUX & Rule & $\times$ & $\checkmark$ & .90 \\
D\textsuperscript{+} & FLUX & Rule+LLM & $\checkmark$ & $\checkmark$ & .90 \\
E\textsuperscript{+} & FLUX & Rule & $\checkmark$ & $\checkmark$ & .90 \\
\midrule
\multicolumn{6}{@{}l}{\emph{v3 VLM scoring (VLM replaces CLIP):}} \\
H & FLUX & VLM & $\times$ & $\checkmark$ & .90 \\
I & FLUX & VLM & $\checkmark$ & $\checkmark$ & .90 \\
J & FLUX & VLM+LLM+Enh & $\checkmark$ & $\checkmark$ & .90 \\
K & FLUX & VLM+LLM & $\checkmark$ & $\checkmark$ & .90 \\
\midrule
\multicolumn{6}{@{}l}{\emph{v4 generator upgrade (NB2 replaces FLUX):}} \\
G & NB2 & VLM & $\times$ & $\checkmark$ & .90 \\
G\textsuperscript{+} & NB2 & VLM & $\times$ & $\checkmark$ & .90 \\
G\textsuperscript{++} & NB2 & VLM & $\checkmark$ & $\checkmark$ & .90 \\
\bottomrule
\end{tabular}
\caption{Ablation conditions. Loop = multi-round (max 3). Route = cultural pipeline routing. $\theta$ = Queen accept threshold. D is the full v1 system; v2 supplementary conditions use tightened $\theta$ to ensure multi-round; v3 VLM conditions replace CLIP scoring with VLM per-layer evaluation; v4 conditions replace FLUX with the NB2 AR+Diffusion hybrid generator. G\textsuperscript{+} injects the full EvidencePack into NB2's 32K context.}
\label{tab:conditions}
\end{table}

Key comparisons:
\begin{itemize}
    \item \textbf{C vs D}: Full pipeline impact (core hypothesis)
    \item \textbf{C vs C\textsuperscript{+}}: Routing-only contribution (no loop, no LLM)
    \item \textbf{C\textsuperscript{+} vs E\textsuperscript{+}}: Loop contribution (controlling for routing)
    \item \textbf{C\textsuperscript{+} vs D\textsuperscript{+}}: Full system vs.\ routing baseline
    \item \textbf{A vs B}: Pipeline impact on weak model
    \item \textbf{A vs C}: Model upgrade effect (SD1.5 $\to$ FLUX)
    \item \textbf{H vs G}: Generator upgrade under VLM scoring (FLUX $\to$ NB2)
    \item \textbf{G vs G\textsuperscript{++}}: Self-correction value on a strong generator
\end{itemize}

\subsection{Evaluation Metrics}

\paragraph{Automatic.} \Lx{1}--\Lx{5} scores (rule-based and LLM hybrid), weighted total score, acceptance rate, rerun trigger rate, round count, and latency.

\paragraph{Human Evaluation.} Formal human evaluation with cultural domain experts is outside the scope of this study; we rely on automatic \Lx{1}--\Lx{5} metrics validated by the blind evaluation pilot (\S\ref{sec:results}).

% ============================================================
\section{Results}
\label{sec:results}
% ============================================================

\begin{figure*}[t]
    \centering
    \includegraphics[width=\textwidth]{figures/qualitative.pdf}
    \caption{Qualitative comparison between single-pass FLUX baseline (C) and the full \system{} pipeline (D) across three representative tasks. The full system achieves substantial improvements in \Lx{5} Philosophical Aesthetics ($+0.256$ to $+0.268$), with the taboo task (left) requiring 3 rounds of self-correction.}
    \label{fig:qualitative}
\end{figure*}

\subsection{Mock Ablation (Framework Validation)}

We first validated framework correctness with deterministic mock scoring (180 runs, 6$\times$30).
Cultural routing is the dominant factor: routed conditions achieve 47--93\% Queen acceptance vs.\ 23\% without, driven by tradition-specific \Lx{5} weight adjustment.
Multi-round looping provides additional uplift (93\% vs.\ 47\% among routed conditions).
Model upgrade shows zero mock effect (expected, as mock scoring is model-independent).

\subsection{Category Breakdown}

\begin{table}[t]
\centering
\small
\begin{tabular}{@{}ccccc@{}}
\toprule
 & \textbf{Poetic} & \textbf{Cultural} & \textbf{Taboo} & \textbf{All} \\
\midrule
A & .824\,{\scriptsize(.016)} & .824\,{\scriptsize(.011)} & .777\,{\scriptsize(.117)} & .809 \\
C & .824\,{\scriptsize(.016)} & .824\,{\scriptsize(.013)} & .781\,{\scriptsize(.121)} & .810 \\
\textbf{D} & \textbf{.880}\,{\scriptsize(.045)} & \textbf{.852}\,{\scriptsize(.016)} & \textbf{.809}\,{\scriptsize(.082)} & \textbf{.847} \\
\midrule
$\Delta$ C{\scriptsize$\to$}D & +.056 & +.028 & +.028 & +.038 \\
\bottomrule
\end{tabular}
\caption{Real-mode category breakdown. SD in parentheses. Poetic tasks show the largest gain ($+0.056$); taboo tasks reduce variance (SD $.121 \to .082$).}
\label{tab:category}
\end{table}

Poetic tasks benefit most ($+0.056$), reflecting the LLM critic's ability to evaluate \Lx{5} philosophical depth; taboo tasks show the largest variance reduction (SD $.121{\to}.082$), suggesting the self-correction loop stabilizes culturally sensitive cases.

\subsection{Blind Evaluation Pilot}

A blind pilot (30 tasks $\times$ 2 groups, 60 real SD1.5 images) validated the hybrid critic: the LLM scores \Lx{5} Philosophical Aesthetics 0.258 points lower than rule-only ($0.629$ vs.\ $0.887$), identifying genuine cultural misalignment that keyword scoring misses.

\subsection{Real-Mode Ablation}
\label{sec:real_ablation}

We ran all 16 conditions in real mode across four phases: SD1.5 via local GPU (RTX 2070, fp16), FLUX.1-Schnell via Together.ai API, and NB2 via Gemini image generation API; v3--v4 VLM scoring uses Gemini~2.5 Flash as the per-layer visual evaluator.
Table~\ref{tab:real_ablation} reports the results.

\begin{table*}[t]
\centering
\small
\begin{tabular}{@{}clcccc|ccccc@{}}
\toprule
 & \textbf{Configuration} & \textbf{Mean} & $\boldsymbol{\Delta}$\textbf{C} & \textbf{$p$} & \textbf{$d$} & \Lx{1} & \Lx{2} & \Lx{3} & \Lx{4} & \Lx{5} \\
\midrule
A & SD1.5 / Rule / Single        & .809 & $-$.001 & .615       & ---  & .846 & .764 & .811 & .933 & .697 \\
B & SD1.5 / Rule+LLM / Multi     & .831 & +.021   & .013*      & 0.48 & .818 & .741 & .838 & .959 & .784 \\
C & FLUX / Rule / Single          & .810 & ---     & ---        & ---  & .821 & .773 & .816 & .933 & .705 \\
\textbf{D} & \textbf{FLUX / Full Pipeline} & \textbf{.847} & \textbf{+.038} & $<$\textbf{.001}*** & \textbf{0.94} & .800 & .766 & \textbf{.872} & \textbf{.964} & \textbf{.839} \\
E & FLUX / Rule / Multi           & .846 & +.037   & $<$.001*** & 1.03 & .816 & .772 & .882 & .933 & .817 \\
F & FLUX / Rule+LLM / Single     & .829 & +.019   & .059       & 0.36 & .799 & .754 & .827 & .964 & .814 \\
\midrule
\multicolumn{6}{@{}l}{\emph{v2 supplementary ($\theta_{\text{accept}} = 0.90$, rerun rate 77\%):}} \\
C\textsuperscript{+} & FLUX / Rule / Single / +Route & .846 & +.036 & $<$.001*** & 0.95 & .814 & .771 & .883 & .933 & .814 \\
D\textsuperscript{+} & FLUX / Full / Multi / +Route  & .837 & +.028 & .002**     & 0.60 & .804 & .759 & .869 & .943 & .818 \\
E\textsuperscript{+} & FLUX / Rule / Multi / +Route  & .848 & +.038 & $<$.001*** & 1.00 & .812 & .773 & .884 & .933 & .824 \\
\midrule
\multicolumn{6}{@{}l}{\emph{v3 VLM scoring (VLM replaces CLIP, $\theta=0.90$):}} \\
H & FLUX / VLM / Single          & \textbf{.887} & \textbf{+.077} & $<$\textbf{.001}*** & \textbf{2.10} & .870 & .848 & .909 & .928 & .871 \\
I & FLUX / VLM / Multi           & .887 & +.077 & $<$.001*** & 1.99 & .875 & .851 & .904 & .929 & .859 \\
J & FLUX / VLM+LLM+Enh / Multi  & .819 & +.009 & .673        & 0.08 & .814 & .752 & .816 & .925 & .816 \\
K & FLUX / VLM+LLM / Multi      & .838 & +.028 & .261        & 0.21 & .816 & .776 & .863 & .922 & .805 \\
\midrule
\multicolumn{6}{@{}l}{\emph{v4 generator upgrade (NB2 AR+Diffusion hybrid):}} \\
G & NB2 / VLM / Single            & .915 & +.105 & $<$.001*** & 3.65 & .896 & .888 & .950 & .930 & .900 \\
G\textsuperscript{+} & NB2 / VLM+Evid / Single   & .908 & +.098 & $<$.001*** & 2.02 & .891 & .876 & .939 & .931 & .892 \\
G\textsuperscript{++} & NB2 / VLM / Multi         & \textbf{.917} & \textbf{+.107} & $<$\textbf{.001}*** & \textbf{2.97} & .899 & .889 & .946 & .930 & .909 \\
\midrule
\multicolumn{2}{@{}l}{$\Delta$ C$\to$D} & & & & & $-$.021 & $-$.008 & \textbf{+.056} & +.030 & \textbf{+.134} \\
\multicolumn{2}{@{}l}{$\Delta$ C$\to$C\textsuperscript{+} (routing)} & & & & & $-$.007 & $-$.002 & \textbf{+.067} & .000 & \textbf{+.109} \\
\multicolumn{2}{@{}l}{$\Delta$ H$\to$G (generator)} & & & & & +.026 & +.040 & \textbf{+.040} & +.002 & +.029 \\
\bottomrule
\end{tabular}
\caption{Real-mode ablation (16$\times$30 = 480 runs) with per-dimension breakdown. v1 conditions (A--F) use $\theta_{\text{accept}}=0.80$; v2 supplementary (C\textsuperscript{+}/D\textsuperscript{+}/E\textsuperscript{+}) use $\theta=0.90$ to ensure multi-round execution; v3 (H/I/J/K) replace CLIP with VLM per-layer scoring; v4 (G/G\textsuperscript{+}/G\textsuperscript{++}) replace FLUX with NB2. $p$-values from paired $t$-tests vs.\ baseline~C; primary results ($p < 0.001$) survive Bonferroni correction for 15 comparisons ($\alpha_{\text{adj}} = 0.003$); secondary comparisons (B: $p = .013$; I$\to$K: $p = .029$) are reported at uncorrected thresholds and should be interpreted cautiously. \textbf{Bold} = peak configuration.}
\label{tab:real_ablation}
\end{table*}

The core comparison C$\to$D yields $+0.038$ ($p < 0.001$, Cohen's $d = 0.94$, a \emph{large} effect by conventional standards).
To decompose this effect, we designed the v2 supplementary experiment with three additional conditions.
The critical finding: \textbf{cultural routing alone} (C$\to$C\textsuperscript{+}) accounts for $+0.036$ ($p < 0.001$, $d = 0.95$)---\textbf{95\% of the full pipeline's effect}.
Adding multi-round iteration on top of routing (C\textsuperscript{+}$\to$E\textsuperscript{+}) yields only $+0.002$ ($p = 0.35$, not significant), despite successfully triggering reruns in 77\% of tasks under the tightened $\theta = 0.90$ threshold.

This result reveals that tradition-specific dimension weighting is the primary mechanism by which \system{} improves cultural quality.
The LLM critic's contribution remains indirect: conditions with LLM scoring (D, F) achieve higher \Lx{4} interpretive scores ($0.964$ vs.\ $0.933$), indicating that the LLM adds depth to cultural interpretation at specific layers.

The per-dimension breakdown (Table~\ref{tab:real_ablation}) confirms that gains concentrate on higher cognitive layers: \Lx{5} ($+0.134$, $p < 0.001$) and \Lx{3} ($+0.056$, $p < 0.001$), while \Lx{1}--\Lx{2} show minor trade-offs ($-0.021$ and $-0.008$, attributable to rule-based scoring sensitivity to longer, culturally enriched prompts).
This confirms that the pipeline specifically improves \emph{cultural understanding}, not generic visual quality (Figure~\ref{fig:qualitative}).

In the v1 experiment ($\theta = 0.80$), most FLUX-based scores exceeded the threshold on the first pass, yielding a low rerun rate (13\% for D).
The v2 supplementary experiment ($\theta = 0.90$) successfully forced multi-round execution (77\% rerun rate, mean 2.5 rounds for D\textsuperscript{+}), yet the additional rounds did not yield significant improvement---revealing that CLIP-based scoring lacks the granularity to differentiate between same-prompt regenerations and thereby cannot drive effective iterative refinement (\S\ref{sec:analysis}).
The A$\to$B comparison confirms pipeline generality across model scales ($+0.022$, $p = 0.008$, $d = 0.52$).

\subsection{Generator Generalization (v4: NB2)}
\label{sec:v4_results}

To test whether the framework generalizes beyond diffusion-only generators, we replaced FLUX with NB2---a hybrid autoregressive-diffusion model with a 32K-token context window and internal ``thinking'' loop.
Three conditions isolate different aspects:
\textbf{G} (NB2 + VLM routing, single pass) establishes the generator baseline;
\textbf{G\textsuperscript{+}} injects the full Scout EvidencePack into NB2's extended context;
\textbf{G\textsuperscript{++}} adds multi-round self-correction (max 3 rounds).

NB2 with VLM routing (\textbf{G}: $0.915$) outperforms the best FLUX configuration (\textbf{H}: $0.887$) by $+0.028$ ($p < 0.001$, $d = 1.01$).
Multi-round iteration (\textbf{G\textsuperscript{++}}: $0.917$) adds only $+0.002$ ($p = 0.70$, not significant)---mirroring the pattern observed with FLUX under VLM scoring (H$\to$I: $+0.000$).
Unexpectedly, injecting the full EvidencePack (\textbf{G\textsuperscript{+}}: $0.908$) slightly \emph{decreases} performance ($-0.007$, $p = 0.31$), with all VLM-scored dimensions trending negative (\Lx{1} $-0.005$, \Lx{2} $-0.012$, \Lx{3} $-0.010$, \Lx{5} $-0.009$) while rule-only \Lx{4} is flat ($+0.001$).
A single taboo task illustrates the mechanism: one task scores $0.915$ under G but drops to $0.725$ under G\textsuperscript{+} ($\Delta = -0.190$), with \Lx{2} collapsing from $0.875$ to $0.525$---consistent with NB2 having already encoded sufficient cultural knowledge internally, and the additional ${\sim}$2K-token EvidencePack introduces attention dilution that degrades technical execution.

Task difficulty is highly consistent across conditions (Pearson $r = 0.86$--$0.95$ for all paired comparisons, $p < 0.001$), confirming that score differences reflect systematic effects rather than random variance.
This pattern---self-correction yielding diminishing returns as generator quality improves---is consistent across all three generator tiers: SD1.5 benefits most from the agent loop ($+0.022$), FLUX benefits from routing but not iteration ($+0.036$ routing, $+0.002$ iteration), and NB2 shows near-zero marginal gain from any pipeline intervention beyond routing.

% ============================================================
\section{Analysis}
\label{sec:analysis}
% ============================================================

\subsection{Why Routing Dominates}

The most striking finding is that cultural routing alone (C$\to$C\textsuperscript{+}: $d = 0.95$) captures 95\% of the full pipeline's effect, while multi-round iteration under the same routing (C\textsuperscript{+}$\to$E\textsuperscript{+}: $d = 0.18$) adds negligible improvement.
This holds even when iteration is \emph{genuinely active}: v2 conditions trigger reruns in 77\% of tasks (mean 2.5 rounds), yet the additional rounds do not yield significant score gains.

We interpret this through two complementary mechanisms.
First, the routing mechanism changes \emph{what is measured}: by amplifying tradition-specific dimensions (e.g., $w_{\Lx{5}} = 0.30$ for xieyi vs.\ $0.20$ uniformly), the router surfaces cultural quality that uniform evaluation suppresses.
The same image scores differently under routed vs.\ unrouted evaluation because the weighting reflects the tradition's actual priorities---this is an evaluation insight, not a generation improvement.
The per-dimension data confirms this: C$\to$C\textsuperscript{+} gains concentrate in \Lx{3} ($+0.067$) and \Lx{5} ($+0.109$), precisely the dimensions amplified by routing.

Second, multi-round iteration fails to improve scores because the underlying scoring mechanism---CLIP-based similarity between the generated image and the cultural prompt---has insufficient granularity to differentiate between same-prompt regenerations.
CLIP scores for different random seeds of the same prompt typically vary by only $\pm 0.01$--$0.02$, well within noise.
The iteration mechanism is architecturally sound (the Queen correctly identifies low-scoring candidates and triggers regeneration), but the scoring signal cannot guide search through generation space when candidate quality variation is below the scorer's discrimination threshold.

\paragraph{VLM Scoring Validation.}
To test whether finer-grained scoring resolves this bottleneck, we ran a supplementary v3 experiment (4 conditions $\times$ 30 tasks; Table~\ref{tab:real_ablation} v3 section) replacing CLIP with VLM per-layer evaluation.
VLM-based routing alone (\textbf{H}: $.887$) outperforms CLIP-based routing (C\textsuperscript{+}: $.846$) by $+0.041$ ($p < 0.001$), confirming evaluation granularity as the binding constraint.
Multi-round iteration under VLM (\textbf{H}$\to$\textbf{I}: $.887 \to .887$, $p = .98$) remains non-significant---VLM resolves the scoring-granularity bottleneck, but reveals a new one: generator output quality already meets the accept threshold, leaving no deficiency for the loop to correct (a ceiling effect confirmed by v4, \S\ref{sec:v4_results}).
Unexpectedly, adding LLM critic alongside VLM (\textbf{I}$\to$\textbf{K}: $.887 \to .838$) \emph{decreases} performance: VLM already subsumes the cultural evaluation function the LLM was designed to provide, producing interference rather than complementarity.
The strongest configuration is FLUX + VLM routing without LLM critic (\textbf{H}: $.887$), outperforming the full CLIP pipeline (\textbf{D}: $.847$) by $+0.040$.
Completing the v3 ablation chain, adding PromptEnhancer to form the full four-factor system (\textbf{K}$\to$\textbf{J}: $.838 \to .819$, $p = 0.48$) shows a consistent negative trend: the Enhancer was optimized for CLIP token coverage and this objective conflicts with VLM evaluation that rewards visual realization directly---demonstrating that component design must co-evolve with the evaluation modality.
Examining K's task-category breakdown reveals that the degradation concentrates in taboo tasks ($\text{K}_{\text{taboo}} = 0.736$ vs.\ $\text{H}_{\text{taboo}} = 0.849$, $\Delta = {-}0.113$): the LLM critic's EvidencePack-based taboo scanning produces conservative scores that conflict with the VLM's direct visual assessment, amplifying interference on the tasks where cultural risk is highest.

\subsection{Why L5 Shows the Largest Gain}

The concentration of improvement in \Lx{5} Philosophical Aesthetics ($+0.134$, $p < 0.001$) reflects the frontier nature of this layer: single-pass models fail most dramatically here (baseline $0.705$) because philosophical resonance requires alignment across visual, cultural, and interpretive dimensions simultaneously.
The routing mechanism directly addresses this by amplifying \Lx{5} for traditions that demand philosophical depth (xieyi: $w_{\Lx{5}} = 0.30$ vs.\ $0.20$ uniformly), which propagates through evidence retrieval and prompt construction to orient generation toward culturally resonant output.

\subsection{Effect Decomposition: Framework vs.\ Generator}

The full improvement from baseline C ($0.810$) to peak G\textsuperscript{++} ($0.917$) decomposes into two orthogonal contributions:
(1)~\textbf{Framework contribution} (C$\to$H: $+0.077$, $d = 2.10$)---replacing CLIP with VLM-based routing, which changes \emph{how} quality is measured and weighted, accounting for 73\% of total improvement; and
(2)~\textbf{Generator contribution} (H$\to$G: $+0.028$, $d = 1.01$)---replacing the FLUX diffusion backend with NB2's hybrid architecture, which changes \emph{what} is generated, accounting for 27\%.
(Note that this decomposes the full path C$\to$G\textsuperscript{++}; the separate finding that routing accounts for 95\% of C$\to$D refers to the v1 CLIP-based experiment scope.)

This 73:27 decomposition reveals that cultural art quality is primarily accessible through \emph{evaluation} improvement: the framework's tradition-specific routing surfaces cultural quality dimensions that uniform evaluation suppresses, and this evaluation insight alone yields larger gains than upgrading the generator.
However, per-dimension analysis shows the contributions are complementary: the framework concentrates gains in \Lx{3}--\Lx{5} (cultural, interpretive, philosophical) while the generator upgrade uniformly raises all dimensions (H$\to$G: \Lx{2} $+0.040$, \Lx{3} $+0.040$, \Lx{5} $+0.029$), with notably higher gains on VLM-scored dimensions (avg $+0.034$) than the rule-only \Lx{4} ($+0.002$, $p = 0.34$)---a pattern warranting further investigation (\S Limitations).

\subsection{Diminishing Returns of Self-Correction}

A consistent pattern emerges across generator tiers: as generator quality increases, the marginal benefit of iterative self-correction decreases.
SD1.5 benefits from the full pipeline ($+0.022$, $p = 0.008$); FLUX benefits from routing ($+0.036$, $p < 0.001$) but not from additional iteration ($+0.002$, $p = 0.35$); NB2 shows near-zero gain from \emph{any} pipeline intervention beyond routing ($+0.002$, $p = 0.70$).
Correlating baseline quality against loop benefit across five paired comparisons yields Pearson $r = -0.81$---a strong negative trend indicating monotonic diminishing returns.

The mechanism is revealed by G\textsuperscript{++}'s round statistics: despite being configured for up to 3 rounds of self-correction, \emph{all 30 tasks were accepted on the first pass} (mean rounds $= 1.0$).
NB2's output quality exceeds the VLM Critic's accept threshold ($\theta = 0.90$) immediately, leaving no deficiency for the loop to correct.
By contrast, FLUX under VLM scoring (condition I) triggers reruns in 40\% of tasks (12/30 multi-round), confirming that the loop architecture functions correctly when the generator leaves room for improvement.
The score distribution further supports a ceiling effect: 90\% of NB2 tasks (27/30) exceed $0.90$, compared to 67\% for FLUX+VLM (H) and 0\% for the CLIP baseline (C).

\subsection{Evaluation-Side vs.\ Generation-Side Correction}

Our approach is fundamentally \emph{evaluation-side}: we wrap an external cognitive loop around an existing generator, diagnosing and correcting cultural deficiencies through an independent evaluation pipeline.
This contrasts with emerging \emph{generation-side} approaches that build self-correction into the generation process itself---for example, models that produce internal ``thought images'' before final output, performing implicit self-evaluation within the same forward pass.

These paradigms are complementary: generation-side correction excels at \Lx{1}--\Lx{2} properties assessable from the visual signal alone, while evaluation-side routing is necessary for \Lx{3}--\Lx{5} because these require \emph{external} cultural knowledge absent from the generator's training distribution---if training data underrepresents Islamic aniconism, internal reasoning is unlikely to recover that convention.
Our results confirm this: \Lx{3} and \Lx{5} show the largest gains ($+0.056$ and $+0.134$) while \Lx{1}--\Lx{2} show minor trade-offs, confirming that evaluation-side routing targets cultural layers that generation-side correction alone does not appear to address.
The v4 NB2 experiment provides direct evidence: NB2 with its internal thinking loop achieves strong \Lx{1}--\Lx{2} scores ($0.896$, $0.888$) through generation-side self-correction, while our framework adds \Lx{3}--\Lx{5} cultural routing that NB2's internal reasoning alone does not appear to replicate.
The full 480-run experiment costs under~\$15 (FLUX API, NB2 API, SD1.5/RTX~2070, Gemini~Flash).

% ============================================================
\section{Conclusion}
\label{sec:conclusion}
% ============================================================

We presented \system{}, a multi-agent evaluation pipeline that addresses cultural decontextualization in AI art generation through tradition-specific routing grounded in the \Lx{1}--\Lx{5} cognitive hierarchy---a synthesis of Panofsky's iconographic analysis and Barrett's criticism pedagogy, operationalized for computational evaluation.
Across 480 real-mode runs (16 conditions $\times$ 30 tasks), VLM-based routing with the NB2 generator achieves the peak score of $\mathbf{0.917}$ ($+0.107$ vs.\ baseline, $d = 2.97$).
Effect decomposition reveals that 73\% of the improvement stems from the evaluation framework (tradition-specific VLM routing) and 27\% from the generator upgrade (FLUX$\to$NB2), confirming that cultural art quality is primarily accessible through evaluation improvement.

Four findings carry implications beyond our specific system.
First, \emph{routing matters more than iteration}: tradition-specific dimension weighting accounts for 95\% of the observed improvement under CLIP scoring ($d = 0.95$), while multi-round iteration adds only $+0.002$ (not significant), identifying scoring granularity as the bottleneck for effective self-correction loops.
Second, \emph{cultural quality is not visual quality}: improvements concentrate in \Lx{3}--\Lx{5} (cultural, interpretive, philosophical) while \Lx{1}--\Lx{2} (visual, technical) show minor trade-offs, confirming that these are distinct axes requiring distinct interventions.
Third, \emph{evaluation modality constrains component design}: VLM-based scoring outperforms CLIP by $+0.041$, but shows substantial functional overlap with LLM critics (I$\to$K: $-0.049$, $p = .029$ uncorrected), and Prompt Enhancers optimized for CLIP token coverage are not beneficial under VLM scoring (K$\to$J: $-0.019$, not significant).
Fourth, \emph{self-correction has diminishing returns}: the agent loop is most valuable for weaker generators (SD1.5: $+0.022$) and progressively less effective for stronger ones (FLUX: $+0.002$ ns; NB2: $+0.002$ ns), suggesting that external evaluation frameworks are most impactful where generation quality leaves room for correction.
Our evaluation-side approach is complementary to generation-side self-correction: NB2's internal thinking loop achieves strong \Lx{1}--\Lx{2} visual quality, while external VLM-based routing provides the independent \Lx{3}--\Lx{5} cultural validation that generators cannot perform on themselves.

% ============================================================
\section*{Limitations}
% ============================================================

\begin{enumerate}
    \item \textbf{Scoring granularity}: CLIP-based scoring cannot differentiate between same-prompt regenerations ($\pm 0.01$--$0.02$ variance), preventing effective multi-round self-correction. Our v3 VLM experiment confirms that VLM-based scoring resolves this limitation ($+0.041$), while revealing a new design constraint: VLM and LLM critic functions overlap under high-quality scoring, requiring careful architectural separation.
    \item \textbf{Potential scorer--generator affinity}: The v4 experiment uses Gemini as both generator (NB2) and evaluator (VLM Critic). Our control dimension \Lx{4} (rule-only, no VLM input) shows near-zero NB2 advantage ($+0.002$, $p = 0.34$), while VLM-scored dimensions show $+0.026$--$+0.040$ (avg $+0.034$, all $p < 0.005$). This asymmetry is consistent with scorer--generator affinity bias~\cite{panickssery2024llm}, though genuine quality differences cannot be excluded. Cross-family evaluation (e.g., Claude as VLM scorer on NB2 output) would provide stronger controls.
    \item \textbf{Human evaluation}: Results rely on automatic \Lx{1}--\Lx{5} metrics; human evaluation with cultural domain experts is needed to validate that score improvements correspond to perceptible quality gains.
    \item \textbf{Generator reproducibility}: NB2 is a closed-source API; our v4 results depend on this specific model's capabilities and may not generalize to other hybrid AR+Diffusion architectures. The framework itself is generator-agnostic by design.
\end{enumerate}

% ============================================================
\section*{Ethics Statement}
% ============================================================

This work generates AI art inspired by diverse cultural traditions. We emphasize that our system is designed to \emph{respect} cultural conventions (including taboos) rather than appropriate or trivialize them.
The taboo detection system specifically prevents generation of culturally offensive content.
All generated images are for research purposes and do not represent authentic cultural artifacts.

% ============================================================
\section*{Data Availability}
% ============================================================

The evaluation tasks are drawn from \vulca{}-Bench~\cite{yu2025vulcabench}, publicly available at \url{https://arxiv.org/abs/2601.07986}.
All 480 ablation run records (16 conditions $\times$ 30 tasks), including per-dimension scores, round counts, and latencies, will be released upon publication.

% ============================================================
% References
% ============================================================
\bibliographystyle{ACM-Reference-Format}
\bibliography{references}

\end{document}
